\section{Preparation of the pair measures}\label{prep:section}

Proposition~\ref{prep:starting-pairs} summarizes the output of this
Section: decreasing pair filters of positive limiting mass with a
controlled product partition. At each scale, a further power-small
deletion gives cap bounds of exponent \(S\) with an arbitrarily small
power loss. Auxiliary projections estimate directions only; all
distances remain those of the original pairs. We first prove the
improvement needed in odd dimensions and then construct the measures
to which it applies.

\subsection{From integer to half-integer directional bounds}

For separated supports, write $\pi_x(z)=(z-x)/|z-x|$. If $G$ is a
set of pairs with pin law $\nu$ and source law $\rho$, its directional
measures are
\[
 (\pi_x)_\#(\rho|_{G_x}),\qquad
 (\pi_z)_\#(\nu|_{G^z}).
\]
These are subprobabilities; they are never normalized separately.
A directional plate means $\{w\in S^k:\dist(w,V)\le u\}$,
where $V$ is a linear hyperplane in $\R^{k+1}$.

\begin{definition}\label{prep:availability}
Let $\rho,\nu$ be fixed separated probabilities in $\R^{k+1}$, and
let $G_N$ be Borel pair sets. An exponent $q$ is \emph{available} for this
family if, for every $e>0$, there are $c_e>0$, $C_e<\infty$, and Borel
$B_N\subset G_N$, for all sufficiently large $N$, such that
\[
 (\nu\otimes\rho)(B_N)\le C_e2^{-c_eN},
\]
and both directional measures of $G_N\setminus B_N$ give every spherical
cap of radius $u\in[2^{-N},1]$ mass at most $2^{eN}u^q$.
The estimates hold at every pin in the retained supports, after removal of
null fibers if necessary.
\end{definition}

\begin{lemma}[A contact interval]\label{prep:contact}
Suppose $m:[0,1]\to\R$ is nondecreasing and Lipschitz,
$m(0)=0$, and $m(u)\ge su$. If $q<t<s$, there exist
$0<A<B<1$, with $B<2A$, such that
\begin{align}
 m(w)&\ge qw &&(0\le w\le A),\label{prep:contact-first}\\
 m(w)-m(A)&\ge t(w-A)&&(A\le w\le B),\label{prep:contact-middle}\\
 m(w)-m(B)&\ge q(w-B)&&(B\le w\le1).\label{prep:contact-last}
\end{align}
\end{lemma}
\begin{proof}
Put $g(u)=m(u)-qu$ and $F(u)=\min_{u\le v\le1}g(v)$.
The function $F$ is nondecreasing and Lipschitz: if $u<v$, a minimizer
for $F(u)$ either lies in $[v,1]$, or comparison with $g(v)$ bounds
$F(v)-F(u)$ by $\Lip(g)(v-u)$. Moreover, $F(0)=0$ and
$F(1)\ge s-q>t-q$. Absolute continuity gives an interior
differentiability point $A$ at which $F'(A)>t-q$.
At this point $F(A)=g(A)$, since a strict gap would force the minimum
to be attained a fixed distance to the right and would make $F$ locally
constant. There are contact points $B>A$ arbitrarily close to $A$:
otherwise $F$ would be constant immediately to the right of $A$.
Choose such a $B<\min(2A,1)$ so close that
$F(w)-F(A)\ge(t-q)(w-A)$ on $[A,B]$.
Using $g(w)\ge F(w)$ proves Equation~\eqref{prep:contact-middle}.
The first inequality follows from $m(w)\ge sw$, and the last follows
from $g(w)\ge F(B)=g(B)$ for $w\ge B$.
\end{proof}

The argument below has precedents in the thin-tube bootstrap of Orponen,
Shmerkin and Wang~\cite[Section~1.3 and Lemma~2.8]{OSW2024} and in
Shmerkin and Wang's use of scale profiles to improve radial
bounds~\cite[Section~5.5]{SW2025}. We use the contact interval to combine
three local estimates while keeping a power bound for the deleted pair mass.

\begin{lemma}[One-scale improvement]\label{prep:one-scale}
Fix $k\ge2$, $k-\tfrac12<s<k$, $k-1\le q<k-\tfrac12$,
and $\zeta>0$. Fix bounded support, separation, and $s$-Frostman
constants, as well as a plate-bound constant $C$. There are $\eps_0,c>0$
such that the following holds for
all sufficiently small dyadic $\delta$. Let $\rho,\nu$ be probabilities
with the fixed bounds, and let $G$ be a Borel pair set whose two
directional measures satisfy, for $\delta\le u\le1$,
\[
 \lambda(B(w,u))\le\delta^{-\eps_0}u^q,
 \qquad
 \lambda\{w:\dist(w,V)\le u\}\le C u^\zeta.
\]
Then one can delete pair mass at most $C\delta^c$ so that both remaining
directional measures give every cap of radius $\delta$ mass at most
$\delta^{q+c}$.
The constants are uniform over the measures and pair sets with these bounds.
\end{lemma}
\begin{proof}
It is enough to prove the conclusion in one orientation and intersect the
two resulting pair sets. In this orientation regularize only the source
law, using Lemma~\ref{reg:regularization}. We first establish the
following uniform assertion for an individual normalized class: sufficiently
small losses in its profile comparisons and in the input cap bound imply
a fixed positive gain and a power bound for the deleted pair mass.
During this class assertion, $\rho$ denotes the normalized class.
Restrict $G$ to pairs whose source belongs to that class, and discard
the null pairs outside $\supp\nu\times\supp\rho$. Deleted pair mass
is measured by $\nu\otimes\rho$.
Only the reverse laws $(\pi_z)_\#(\nu|_{G^z})$ are used below.
They remain unnormalized restrictions of the original pin law, so
normalizing a source class does not alter their cap or plate constants.

We prove this assertion by contradiction, which also fixes its quantifiers.
If it failed, there would be scales $R=\delta^{-1}=2^M\to\infty$,
class-comparison losses tending to zero, input cap losses tending to zero,
and counterexamples to every gain and deletion exponent tending to zero.
The class profiles are equi-Lipschitz. Pass to a uniformly convergent
subsequence with limit $m$, where $m(0)=0$ and $m(u)\ge su$,
by the lower profile bound in Lemma~\ref{reg:regularization} and
the vanishing regularization losses. All normalized cell masses and conditional cell comparisons are
then $R^{o(1)}$ times those given by $m$.

Fix $t\in(k-\tfrac12,s)$ and choose $A,B$ by
Lemma~\ref{prep:contact}. Use the four depths
$0,a=\lfloor AM\rfloor,b=\lfloor BM\rfloor,M$.
Rounding changes only $R^{o(1)}$ factors. We establish local radial
fiber estimates on the three intervals. For an occupied source cell $Q$ at
the initial depth of an interval $[u,v]$, write
$\rho_Q=\rho|_Q/\rho(Q)$ and define
\[
 \mathcal R_{Q,v}(x,z)
 =\rho_Q\{z'\in Q:
       |\pi_x(z')-\pi_x(z)|\le C_{\mathrm{rad}}R^{-v}\},
 \qquad x\in\supp\nu,\quad z\in Q\cap\supp\rho.
\]
Here $C_{\mathrm{rad}}$ is a sufficiently large fixed width constant.
The comparison source $z'$ always ranges over the full conditional
class law in $Q$. We will bound this quantity only at retained witness
pairs $(x,z)$; no individual pair fiber is normalized. The function
$\mathcal R_{Q,v}$ is Borel by parameterized integration.

First consider $[u,v]=[0,A]$ or $[B,1]$, and a source cell $Q$ at
depth $uM$.
For every $u\le w\le v$, the profile comparisons and
Equations~\eqref{prep:contact-first} and~\eqref{prep:contact-last} give
\begin{equation}\label{prep:local-ball}
 \rho_Q(B(z,CR^{-w}))\le R^{-q(w-u)+o(1)}.
\end{equation}
At a fixed witness $z\in Q$, integrate pins only on $G^z$.
Bounded separation implies
\[
 |\pi_x(z)-\pi_x(z')|\le C_{\mathrm{rad}}R^{-v}
 \quad\Longrightarrow\quad
 \pi_z(x)\text{ lies in two caps of radius }
 C'R^{-v}/|z-z'|.
\]
For $\tfrac12R^{-w}<|z-z'|\le R^{-w}$, the product of the cap bound
and Equation~\eqref{prep:local-ball} is at most
$R^{-q(v-u)+o(1)}$; equivalently one may sum over all intermediate
dyadic distance annuli. The innermost ball $|z-z'|\lesssim R^{-v}$
has the same bound by Equation~\eqref{prep:local-ball}. Hence
\begin{equation}\label{prep:local-mean}
 \int_{G^z}\mathcal R_{Q,v}(x,z)\,d\nu(x)
 \le R^{-q(v-u)+o(1)}.
\end{equation}
Fubini and Markov therefore delete only $R^{-e'+o(1)}$ pair mass,
for any fixed $e'>0$, and leave
\[
 \mathcal R_{Q,v}(x,z)\le R^{-q(v-u)+e'}
\]
at every retained witness $(x,z)$ in this cell. Summing over cells
uses the cell weights $\rho(Q)$.

For the middle interval put $D=R^{-(B-A)}$. After rescaling an
$a$-cell $Q$ to bounded size, its conditional probability has a
$t$-ball bound down to scale $D$, with loss $D^{-o(1)}$, by
Equation~\eqref{prep:contact-middle}. For the unchanged reverse laws,
$B-A>0$ is fixed along this subsequence, so their $R^{o(1)}$ cap loss
is $D^{-o(1)}$ on the required radii. At $q=k-1$, their fixed plate
constant also fits the permitted small $D$-power loss for small $D$.
Apply Corollary~\ref{ang:bad-directions} with a sufficiently large fixed
tube-width constant. It supplies a fixed $c_1>0$ and a common Borel
set $\mathcal B_Q$ containing all directions for which a union of
at most $D^{-q-c_1}$ parallel tubes of width $C'D$ carries at least
$D^{c_1}$ conditional source mass. For every $z\in Q$,
\begin{equation}\label{prep:witness-bad}
 ((\pi_z)_\#(\nu|_{G^z}))(\mathcal B_Q)\le D^{c_1}.
\end{equation}
The directional law in this statement is a subprobability and has not
been renormalized. The set $\mathcal B_Q$ depends on $Q$ and its
conditional source law, not on $z$. If necessary decrease $c_1$ so
that the same exponent bounds the exceptional mass. Constant-width
enlargements and gridding are part of this Corollary.

Delete the pairs in Equation~\eqref{prep:witness-bad}, at cost
$O(D^{c_1})$ by Fubini. For a pin still represented in $Q$, at least
one retained witness has a good direction. Let $c_Q$ be the cell center
and $w_{x,Q}=\pi_x(c_Q)$. All witness directions differ, up to sign,
from $w_{x,Q}$ by $O(R^{-A})$. Since
\begin{equation}\label{prep:linearization-margin}
 R^{-A}=o(D),\qquad R^{-2A}=o(R^{-B}),
\end{equation}
a tube cover at $w_{x,Q}$ with a slightly smaller fixed width would
also be a forbidden cover at that witness direction. Thus the
conditional source law is robust at the center direction, uniformly
over every represented pin. The directions may be regarded as
unoriented for this assertion.

For $z,z'\in Q$, Taylor expansion of $y\mapsto\pi_x(y)$ on the
separated supports gives
\[
 |\pi_x(z)-\pi_x(z')|\le C_{\mathrm{rad}}R^{-B}
 \quad\Longrightarrow\quad
 |P_{w_{x,Q}^{\perp}}(z-z')|
 \le C''(R^{-B}+R^{-2A})\le C'''R^{-B}.
\]
Grid the perpendicular projection at this scale. For each pin $x$, let
\[
 \mathcal H_{x,Q}
 =\{z\in Q\cap\supp\rho:\mathcal R_{Q,B}(x,z)>D^{q+c_1/2}\}.
\]
If $z\in\mathcal H_{x,Q}$, some one of a bounded number of neighboring
projected grid cells has full $\rho_Q$-mass at least a fixed multiple
of $D^{q+c_1/2}$. There are at most $C D^{-q-c_1/2}$ such heavy
cells. Enlarge each by its bounded number of neighboring grid cells.
Thus $\mathcal H_{x,Q}$ lies in at most that many tubes, of a fixed
width permitted by the preceding robustness assertion. For small $D$
this is less than $D^{-q-c_1}$. For every represented pin, robustness
therefore gives $\rho_Q(\mathcal H_{x,Q})\le D^{c_1}$.

From the pair relation surviving the bad-direction deletion, delete
the pairs $(x,z)$ with $z\in\mathcal H_{x,Q}$. If a pin has no surviving
witness in $Q$, this deletion has zero fiber mass. Otherwise the preceding
bound applies. Integrating over the full pin law, and then summing
with the cell weights $\rho(Q)$, costs $O(D^{c_1})$. This deletion
intersects the first surviving Borel relation with the threshold event
for $\mathcal R_{Q,B}$; it is Borel and needs no measurable choice of
the witness used in the robustness argument. Every remaining witness
satisfies $\mathcal R_{Q,B}(x,z)\le D^{q+c_1/2}$.

Denote the three local factors, in interval order, by
\[
 f_0=R^{-qA+e'},\qquad
 f_1=R^{-(q+c_1/2)(B-A)},\qquad
 f_2=R^{-q(1-B)+e'}.
\]
Let $G'$ be the common Borel pair relation left in the class-restricted
$G$ after the two outer Markov deletions and the two middle deletions,
before endpoint density pruning. At each $(x,z)\in G'$ in the initial
cell $Q$ of the corresponding interval $[u,v]$, the local estimate is
\[
 \rho\{z'\in Q:
       |\pi_x(z')-\pi_x(z)|\le C_{\mathrm{rad}}R^{-v}\}
 =\rho(Q)\mathcal R_{Q,v}(x,z)\le f_j\rho(Q).
\]
Thus Lemma~\ref{reg:pruning-multiplication} applies at the four
endpoint depths with this $G'$ and with
$E_0=E_1=E_2=\supp\nu\times\supp\rho$: the comparison sources are the
full class law, while the estimates are required only at witnesses in
$G'$. The fixed choice of $C_{\mathrm{rad}}$ includes the enlargement
from rounding the endpoint depths. Use density threshold $R^{-e''}$.
The Lemma deletes
$O(R^{-e''})$ pair mass and multiplies the factors with at most
$R^{3e''}$ loss; its density test concerns $G'$ before this last deletion.
Choose $e',e''>0$ so that
\[
 2e'+3e''<\tfrac14c_1(B-A).
\]
The resulting cap mass is at most
$R^{-q-c_1(B-A)/4}$, and every deletion has a fixed positive
power bound. This contradicts the selected counterexamples. In particular,
although $A,B$ arose from a limiting profile, the contradiction proves
uniform positive constants for sufficiently small fixed comparison losses.

Choose the regularization parameter small enough for this uniform class
assertion, and then choose $\eps_0$ small enough for its input cap
loss. Discarded classes have power-small total mass. Apply the assertion
to every remaining class and recombine with its original source weight.
The weights sum to at most one, so both cap estimates and deletion
estimates sum without a factor equal to the number of classes.
Finally repeat with the two original measures interchanged and intersect
the outputs. Decreasing $c$ absorbs fixed constants and proves the Lemma.
\end{proof}

\begin{proposition}[Thin-tube upgrade]\label{prep:thin-upgrade}
Let $k\ge2$ and $k-\tfrac12<s<k$. Let $\rho,\nu$ be bounded,
separated $s$-Frostman probabilities in $\R^{k+1}$. Suppose that the
two directional measures of $G_N$ have plate bounds $Cu^\zeta$,
uniformly for $2^{-N}\le u\le1$, with fixed $\zeta>0$.
If $k-1$ is available, then $k-\tfrac12$ is available.
\end{proposition}
\begin{proof}
Availability is downward closed. It is also closed under increasing
limits inside $[k-1,k-\tfrac12]$: for a desired loss $e$, choose an
available $q$ within $e/2$ of the supremum, apply its estimate with
loss $e/2$, and use $u\ge2^{-N}$.

Suppose that an available exponent $q<k-\tfrac12$ is given. Let
$c,\eps_0$ be the constants in Lemma~\ref{prep:one-scale}, and choose
$q'>q$ with $q'\le q+c/2$ and $q'\le k-\tfrac12$. For any desired
loss $e>0$, choose $e_1>0$ so small that $2e_1q'<e$.
First use availability of $q$ with input loss
$\eps<e_1\eps_0/2$. For every integer $M\in[e_1N,N]$, that
input loss is at most $2^{\eps_0M}$, so
Lemma~\ref{prep:one-scale} applies at $\delta=2^{-M}$ to the same
remaining pair set. Intersect the outputs over these $M$.
Their combined deletion is at most
\[
 C\sum_{M\ge e_1N}2^{-cM}=O(2^{-ce_1N}),
\]
in addition to the original power-small deletion. At these radii the
output is stronger than the requested $2^{eN}u^{q'}$ bound. At
larger radii that bound is trivial from mass at most one, for all large
$N$, by the choice of $e_1$. Dyadic enlargement is absorbed by the
remaining loss. Hence $q'$ is available. Applying this to the
supremum of the available exponents rules out a supremum below
$k-\tfrac12$; closure gives the endpoint itself.
\end{proof}

\subsection{Starting measures and radial densities}

We record explicitly the classical exceptional case, which allows the
remaining construction to assume that no relevant affine plane is charged.

\begin{lemma}[The classical distance threshold]\label{prep:classical-plane}
Let $\rho$ be a compactly supported probability in $\R^n$, $n\ge2$,
with an $s$-ball bound for some $s>(n+1)/2$. Then the distances
between points of $\supp\rho$ form a set of positive Lebesgue measure.
\end{lemma}
\begin{proof}
Choose positive normalized restrictions $\rho_1,\rho_2$ to sufficiently
small separated compact patches. Put $a=(n-1)/2$ and
\[
 \tau=D_\#\bigl[D^{-a}(\rho_1\otimes\rho_2)\bigr],
 \qquad D(x,y)=|x-y|.
\]
This is a nonzero finite positive measure on the distance line.
The elementary sphere-kernel estimate
\[
 |\widehat{\sigma_{S^{n-1}}}(\xi)|\lesssim(1+|\xi|)^{-a}
\]
follows by a partition into the two stationary patches and their
complement, using stationary phase on the former and integration by
parts on the latter. Consequently
\begin{equation}\label{prep:classical-spherical}
 \int_{S^{n-1}}|\widehat{\rho_1}(rw)|^2\,d\sigma(w)
 \lesssim r^{-a}I_a(\rho_1),\qquad r\ge1.
\end{equation}
The energy is finite since $a<s$. The ball bound and a smooth Fourier
ball majorant give
\begin{equation}\label{prep:classical-annulus}
 \int_{r\asymp R}\int_{S^{n-1}}
       |\widehat{\rho_2}(rw)|^2\,d\sigma(w)\,\frac{dr}{R}
 \lesssim R^{-n}\int_{|\xi|\lesssim R}|\widehat{\rho_2}(\xi)|^2d\xi
 \lesssim R^{-s}.
\end{equation}
For the last inequality, the inverse transform of the majorant is
$O_L(R^n(1+R|x-y|)^{-L})$; integrating by dyadic annuli against
the $s$-ball bound proves the claim.

Choose a smooth angular cutoff equal to one on all the cross directions
and zero near their antipodes. The leading spherical stationary-phase
term, uniformly on the separated patches, gives, up to the harmless
Fourier-sign choice,
\[
 \widehat\tau(r)=c_n^{-1}r^a
 \int_{S^{n-1}}\psi(w)\widehat{\rho_1}(rw)
                 \overline{\widehat{\rho_2}(rw)}\,d\sigma(w)
 +O(r^{-1}),\qquad c_n\ne0.
\]
Indeed the error before multiplication by $r^a$ is $O(r^{-a-1})$,
because distances stay in a compact subset of $(0,\infty)$.
Cauchy--Schwarz, Equations~\eqref{prep:classical-spherical}
and~\eqref{prep:classical-annulus}, and $2a+1=n$ imply
\[
 \int_{r\asymp R}|\widehat\tau(r)|^2\,dr
 \lesssim R^{n-a-s}+R^{-1}
 =R^{(n+1)/2-s}+R^{-1}.
\]
The dyadic sum converges. The negative frequencies are conjugates,
and low frequencies are bounded. Plancherel therefore gives an $L^2$
density for $\tau$, proving the conclusion. This is the classical
threshold argument of Falconer~\cite{Falconer1985}.
\end{proof}

\begin{lemma}[Simultaneous projected Frostman bounds]\label{prep:projections}
Let $E\subset\R^d$ be compact and $\dim_HE>S=d/2$.
There is a probability on $E$ with a ball exponent strictly greater
than $S$. Finitely many successive generic orthogonal projections to
dimensions strictly greater than $S$ preserve finite energies with
exponents strictly greater than $S$. After a positive compact restriction
of the original probability, all these projected laws, as well as the
original law, have a common ball exponent $S+\chi$, with $\chi>0$.

If $d=2k-1$ and the original measure gives zero mass to every affine
$k$-plane, the projection to $\R^{k+1}$ may additionally be chosen
to give zero mass to every affine hyperplane.
\end{lemma}
\begin{proof}
Choose $S<s_0<\dim_HE$ and use Frostman's Lemma to obtain an
$s_0$-ball probability. Fix $u>S$ below $s_0$ and all the target
dimensions. Lemma~\ref{conv:energy} gives finite $u$-energy for
almost every successive projection. To make all the projected ball
bounds hold on the same restriction, use the potential proof of that
Lemma: for each projected law $\eta$, retain a set on which
$U_u\eta$ is bounded, and pull this set back to the original space.
Each discarded preimage can have arbitrarily small mass, so the
intersection of these finitely many restrictions has positive mass.
An inner compact restriction has the desired ball bound in every
space. Decrease the exponent slightly if necessary to obtain one
common $S+\chi$.

For the last assertion, almost every $(k+2)$-tuple of original points
is affinely independent: sample successively and use zero mass on
every span of dimension at most $k$. For each fixed such tuple,
almost every projection to $\R^{k+1}$ preserves its independence.
Fubini makes this true for almost every tuple for almost every
projection, simultaneously with energy preservation. If a projected
hyperplane had mass $a>0$, the dependent tuples would have mass at
least $a^{k+2}$, a contradiction. Subsequent restrictions preserve
this property.
\end{proof}

The following lemma adapts the finite-energy proof of
Orponen~\cite[Section~3 and the proof of Theorem~1.13]{Orponen2019}.

\begin{lemma}[Radial $L^p$ estimate]\label{prep:radial-lp}
Let $h\ge1$ and $0<\epsilon<\min(h,1)$. Let $\alpha,\beta$ be
bounded, separated probabilities in $\R^{h+1}$ with finite
$I_{h+\epsilon}$ energies. For some $p>1$, the radial measures
$(\pi_x)_\#\alpha$ have jointly measurable densities $F(x,w)$ and
\[
 \int\!\int_{S^h}F(x,w)^p\,d\sigma(w)\,d\beta(x)<\infty.
\]
The same assertion holds in the reverse orientation. Positive compact
restrictions can be chosen on both sides so that the two radial
$L^p$ norms are uniformly bounded at every remaining pin.
\end{lemma}
\begin{proof}
Put
\[
 u=h-\epsilon,\qquad p=1+\frac{\epsilon}{4h},
 \qquad t=h-\frac{\epsilon}{2}.
\]
Then $1<p<2$ and $up<t<h(2-p)$. Initially suppose that the
source is smooth. For $w\in S^h$, let $P_w$ be projection to
$w^\perp$, let $f_w$ be the density of $(P_w)_\#\alpha$, and put
$\eta_w=(P_w)_\#\beta$. The radial density is
\[
 F(x,w)=\int_0^\infty\alpha(x+rw)r^h\,dr
 \lesssim f_w(P_wx),
\]
because the supports lie in a fixed bounded region.

If $g\ge0$ and $\|g\|_{L^{p/(p-1)}(\eta_w)}=1$, Hölder on the
product space gives
\[
 I_u(g\eta_w)\le I_{up}(\eta_w)^{1/p}
 \lesssim I_t(\eta_w)^{1/p}.
\]
The second inequality uses bounded support and $up<t$.
The energy identity
$I_u(\eta)=c_{h,u}\int|\widehat\eta(\xi)|^2|\xi|^{u-h}\,d\xi$
follows by writing the Riesz kernel as an integral of Gaussian
kernels; see \cite[Lemma~12.12]{Mattila1995} for this identity with
our Fourier convention. Fourier duality and Cauchy--Schwarz with reciprocal
weights therefore imply
\begin{equation}\label{prep:radial-duality}
 \|f_w\|_{L^p(\eta_w)}
 \lesssim I_t(\eta_w)^{1/(2p)}
 \left(\int_{w^\perp}|\widehat\alpha(\xi)|^2
                       |\xi|^{h-u}\,d\xi\right)^{1/2}.
\end{equation}
The pairing can first be justified with bounded smooth truncations,
then passed to the limit by the same energy bound.

Write $A(w)$ for the Fourier integral in
Equation~\eqref{prep:radial-duality} and $B(w)=I_t(\eta_w)$.
Polar coordinates and rotational invariance give
\[
 \int_{S^h}A(w)\,d\sigma(w)
 =C\int_{\R^{h+1}}|\widehat\alpha(\xi)|^2
                      |\xi|^{h-u-1}\,d\xi
 =C'I_{h+\epsilon}(\alpha).
\]
Set $r_0=1/(2-p)$. The one-vector estimate
$\sigma\{w:|P_wz|\le a|z|\}\lesssim a^h$, together with
$tr_0<h$, shows that
\[
 \big\||P_wz|^{-t}\big\|_{L^{r_0}(d\sigma)}\lesssim|z|^{-t}.
\]
Minkowski gives $\|B\|_{L^{r_0}}\lesssim I_t(\beta)$.
Raise Equation~\eqref{prep:radial-duality} to the $p$th power and
apply Hölder with exponents $2/p$ and $2/(2-p)$ to obtain
\begin{equation}\label{prep:radial-integrated}
 \int\!\int F(x,w)^p\,d\sigma(w)\,d\beta(x)
 \lesssim I_{h+\epsilon}(\alpha)^{p/2}I_t(\beta)^{1/2}.
\end{equation}

For a general source, let $\alpha_\varepsilon$ be its convolution with
a nonnegative smooth approximate identity of scale $\varepsilon$ and
sufficiently small support. Separation persists, and the
Fourier energy formula gives a uniform energy bound. The joint
measures
\[
 \Lambda_\varepsilon
 =\bigl[(x,y)\mapsto(x,\pi_x(y))\bigr]_\#
       (\beta\otimes\alpha_\varepsilon)
\]
have uniformly bounded densities in $L^p(\beta\otimes\sigma)$ by
Equation~\eqref{prep:radial-integrated}. Separation gives weak
convergence to the analogous joint measure $\Lambda$.
Duality against continuous functions, followed by $L^{p'}$ density
\cite[Theorem~6.15 and Proposition~7.9]{Folland1999},
shows that $\Lambda$ has an $L^p$ density $F$. Disintegration,
checked on a countable determining family of continuous functions
on $S^h$, identifies $F(x,\cdot)$ as the radial density for
$\beta$-almost every $x$. A Borel version supplies joint measurability.

Repeat with $\alpha$ and $\beta$ interchanged. Choose positive
compact subsets on which the relevant radial norms are bounded and
on which the density identities hold. Restricting the sources only
decreases positive radial measures; normalization costs a fixed
constant. The two uniform bounds therefore hold simultaneously.
For each resulting radial kernel $K_x$, choose the density version
\[
 \widetilde F(x,w)=\limsup_{n\to\infty}
 \frac{K_x(B(w,2^{-n}))}{\sigma(B(w,2^{-n}))}.
\]
This is jointly Borel by parameterized integration. At every retained
pin the preceding domination gives $K_x\ll\sigma$, so spherical
differentiation identifies this version as its density almost
everywhere in $w$, with the stated $L^p$ bound. Here one may view both
measures as ambient Radon measures and apply
\cite[Chapter~1, Theorems~3.23--3.24]{Simon2014}.
The compact good-pin
sets contain the supports of the restricted laws; hence this
construction introduces no new exceptional pins. Apply the same
choice in the reverse orientation.
\end{proof}

\begin{lemma}[Integer caps and lifting]\label{prep:integer-caps}
Under the uniform conclusion of Lemma~\ref{prep:radial-lp}, for
every $e>0$ and $R=2^N$ one can remove pair mass
$O(R^{-e(p-1)/2})$ so that the two remaining directional measures
give every cap of radius $u$ mass at most $R^eu^h$, for all
$R^{-1}\le u\le1$ and all sufficiently large $R$.

The same cap bound lifts through a fixed orthogonal projection from
a higher-dimensional space, provided the original and projected
cross supports are bounded and separated. Constants may change,
and may be absorbed by a smaller preliminary choice of $e$.
\end{lemma}
\begin{proof}
Use the jointly Borel densities $F(x,w)$ and $H(y,w)$ of the two
radial laws. With $T=R^{e/2}$, delete pairs for which
\[
 F(x,\pi_x(y))>T\quad\text{or}\quad H(y,\pi_y(x))>T.
\]
Their product mass is $O(T^{1-p})$, since
$\int_{F>T}F\le T^{1-p}\int F^p$, and similarly for $H$.
Each remaining directional density is bounded by $T$, so the
spherical cap-volume bound gives $CTu^h\le R^eu^h$ for large $R$.

Let $P$ be the fixed projection. On all actual cross directions,
$|Pw|\ge c>0$, because projected separation is bounded below and
original distances are bounded above. The map
\[
 \Phi(w)=Pw/|Pw|
\]
is Lipschitz on this region and satisfies
$\Phi(\pi_x(y))=\pi_{Px}(Py)$. The image of any original cap,
intersected with the actual directions, is contained in a projected
cap of comparable radius. Pushforward of the source measure gives
the asserted bound at each original pin. This uses no injectivity
of $P$, and works in both orientations.
\end{proof}

\subsection{A decreasing coarse filter in odd dimensions}

For an affine hyperplane $H$, write $H_r=\{z:\dist(z,H)\le r\}$.
The following construction gives the spatial regularity that will survive
all the way to the limiting distance measure.

The proof adapts the heavy-plate localization and exclusion argument of
Shmerkin~\cite[Appendix~B]{Shmerkin2023}, in the quantitative form
developed by Ren~\cite[Section~7]{Ren2023}. Here the exclusions are
unions of coarse cell pairs at successive dyadic scales, producing
decreasing pair filters with positive limiting product mass.

\begin{lemma}[Coarse exclusion of heavy plates]\label{prep:coarse-filter}
Let $k\ge2$, and let $\rho_1,\rho_2$ be bounded separated
$s$-Frostman probabilities in $\R^{k+1}$, with
$k-\tfrac12<s<k$. Suppose $\rho=\rho_1+\rho_2$ gives zero mass
to every affine hyperplane. For every fixed $0<\omega\le1$, there
are decreasing Borel sets $\Gamma_N\subset\supp\rho_1\times\supp\rho_2$
such that
\begin{enumerate}
\item $(\rho_1\otimes\rho_2)(\bigcap_N\Gamma_N)>0$;
\item $\Gamma_N$ is a union of pairs of cells in a fixed dyadic
grid of depth at most $\omega N+O(1)$;
\item both directional measures of $\Gamma_N$ satisfy
$\lambda\{w:\dist(w,V)\le u\}\le C_\omega u^\zeta$
for $2^{-N}\le u\le1$, with some $\zeta>0$.
\end{enumerate}
\end{lemma}
\begin{proof}
Put $a=s-(k-1)>1/2$. Covering a bounded portion of an affine
$(k-1)$-plane by $O(r^{-(k-1)})$ balls gives
\begin{equation}\label{prep:lower-plane}
 \rho(L_r)\le C r^a
\end{equation}
uniformly over such planes and $0<r\le1$. This also bounds
neighborhoods of all lower-dimensional affine subspaces.
We shall fix a sufficiently small $b>0$ and then a sufficiently large
integer $j_0$.

Call $H$ heavy at width $r$ if $\rho(H_r)\ge r^b$.
We first prove that its parameter set has an $r$-net of size
\begin{equation}\label{prep:heavy-net}
 O(r^{-C_0b}),
\end{equation}
where $C_0$ is independent of sufficiently small $b$.
Restrict to the compact family of hyperplanes meeting a fixed bounded
neighborhood of the supports, and use the usual metric given by
normal direction, modulo sign, and offset. Fix $L>2/a$.
Within $H_r$, choose $k+1$ points successively, each at distance at
least $r^{Lb}$ from the affine span of its predecessors.
Equation~\eqref{prep:lower-plane} bounds each forbidden set by
$Cr^{Lab}\le r^b/2$ when $r$ is sufficiently small. Thus the
product measure of admissible tuples is at least $c r^{(k+1)b}$,
and every such tuple has $k$-simplex volume at least $c_k r^{kLb}$.

If a fixed admissible tuple belongs to the width-$r$ plates of two
hyperplanes, elementary linear algebra bounds the parameter distance
between these hyperplanes by $Cr^{1-kLb}$. Indeed the matrix of
the $k$ difference vectors has bounded operator norm and smallest
nonzero singular value at least a constant times its $k$-volume;
both normals have scalar products $O(r)$ with every difference.
The offset is then controlled by either first vertex. The hyperplane
parameter space has dimension $k+1$, so an $r$-separated list
contains at most $Cr^{-k(k+1)Lb}$ hyperplanes compatible with this
tuple. Integrating tuple incidences over any $r$-separated list of
heavy hyperplanes proves Equation~\eqref{prep:heavy-net}, with,
for example, $C_0=(k+1)(1+kL)$ after increasing fixed constants.
A maximal separated list is the desired net. Its plates enlarged to
width $C_kr$ contain every original width-$r$ heavy plate in the
bounded region, as well as every grid cell of side $r$ meeting one.

For each $j\ge j_0$, put $r_j=2^{-j}$ and delete every pair of
depth-$j$ cells which both meet some heavy width-$r_j$ plate.
These are finite unions of products of Borel cells. Let $E_j$ be
the deleted pair set. With $J_N=\lfloor\omega N\rfloor$, define
\begin{equation}\label{prep:gamma-definition}
 \Gamma_N=(\supp\rho_1\times\supp\rho_2)
              \setminus\bigcup_{j_0\le j\le J_N}E_j.
\end{equation}
The sets decrease, and nested grids give the stated cell partition.
We prove positive limiting mass by bounding new losses over all $j$.

At a fixed width $r$, use the enlarged plates of its net. Two such
plates with normal angle greater than $\sqrt r$, if they intersect
the bounded region, have their intersection within an
$O(\sqrt r)$-neighborhood of an affine $(k-1)$-plane.
This follows by solving the two normal equations: the smallest
singular value of the pair of unit normals is comparable to their
angle. Equations~\eqref{prep:lower-plane} and~\eqref{prep:heavy-net}
therefore show that the union of all these intersections has
$\rho_1$-mass at most
\begin{equation}\label{prep:plate-exception}
 C r^{a/2-2C_0b}.
\end{equation}
Outside this exceptional set, all listed plates containing a fixed
first endpoint have pairwise normal angle at most $\sqrt r$.
Their bounded portions lie in one plate of width $C\sqrt r$:
the common endpoint controls the offsets as well as the normals.
Consequently all second endpoints lost at this scale lie in that
single plate.

For $j>4j_0$, set $j'=\lfloor j/3\rfloor$.
Taking $j_0$ large ensures $j'\ge j_0$ and
$C2^{-j/2}<2^{-j'}$. The single plate just obtained is contained
in a width-$2^{-j'}$ plate containing the first endpoint.
If this enclosing plate is heavy at that earlier width, every pair
under consideration was already deleted by $E_{j'}$: it contains
the actual endpoints, so both of their depth-$j'$ cells meet it.
If it is light, all newly affected second endpoints have total
$\rho_2$-mass at most $2^{-bj'}$. It follows from
Equation~\eqref{prep:plate-exception} that the new product mass
lost at scale $j$ is at most
\begin{equation}\label{prep:late-loss}
 C2^{-j(a/2-2C_0b)}+2^{-b\lfloor j/3\rfloor}.
\end{equation}

Treat the initial scales $j_0\le j\le4j_0$ together.
Their combined net has at most $C_b2^{4C_0bj_0}$ members.
Enlarge them all to width $C2^{-j_0}$ and repeat the intersection
argument at angle $2^{-j_0/2}$. The exceptional first-endpoint
mass is at most
\[
 C_b2^{-(a/2-8C_0b)j_0}.
\]
For every other first endpoint, all lost second endpoints lie in
one width-$C2^{-j_0/2}$ hyperplane plate.
The hypothesis of zero hyperplane masses implies
\begin{equation}\label{prep:uniform-vanishing}
 \sup_H\rho(H_\varepsilon)\longrightarrow0
       \qquad(\varepsilon\downarrow0).
\end{equation}
Indeed a contrary sequence of hyperplanes has a convergent subsequence
in the compact parameter space; its shrinking plates would force
positive mass on the limiting hyperplane. Thus the total initial
loss tends to zero as $j_0\to\infty$.

Choose $b>0$ so small that $8C_0b<a/4$ and $kLb<1/2$.
The sum of Equation~\eqref{prep:late-loss} over $j>4j_0$ also
tends to zero as $j_0\to\infty$. Choose $j_0$ so large that the
combined loss is less than $1/2$. This proves positive limiting mass.

Finally fix a pin $x$. Directions in a linear-hyperplane
$u$-neighborhood correspond to sources in a width-$Cu$ affine
plate through $x$. At an available comparable exclusion scale,
a heavy plate leaves no retained pairs, while a light plate has
source mass at most its width to the power $b$. At radii below the
last exclusion scale use the last scale instead. The finitely many
radii above the first scale cost a fixed constant. Thus both fibers
obey
\[
 \lambda\{w:\dist(w,V)\le u\}
 \le C\max(u,2^{-J_N})^b
 \le C_\omega u^{\omega b},
 \qquad 2^{-N}\le u\le1.
\]
Dyadic rounding and fixed width constants change only $C_\omega$.
We may take $\zeta=\omega b$, or any smaller positive exponent.
\end{proof}

\subsection{Starting pair measures}

\begin{proposition}[Starting pair measures]\label{prep:starting-pairs}
Let $E\subset\R^d$ be compact, $d\ge2$, and
$\dim_HE>S=d/2$. Either Lemma~\ref{prep:classical-plane} already
proves that its distance set has positive Lebesgue measure, or the
following construction is possible.

After a similarity there are compact separated sets $E_1,E_2\subset E$
and probabilities $\mu_i$ supported on them, each satisfying
\[
 \mu_i(B(x,u))\le C u^{S+\chi}
\]
for a fixed $\chi>0$. Their cross directions lie in a fixed
arbitrarily small spherical patch. For every sufficiently small fixed
$\omega>0$, there are decreasing Borel pair sets
$\Gamma_N\subset E_1\times E_2$ such that
\begin{enumerate}
\item $(\mu_1\otimes\mu_2)(\bigcap_N\Gamma_N)>0$;
\item $\Gamma_N$ is a disjoint union of at most $C2^{2d\omega N}$
products of spatial Borel sets;
\item for every $e>0$, there are $c(e)>0$, $C_e<\infty$ and
Borel $B_N\subset\Gamma_N$ with
$(\mu_1\otimes\mu_2)(B_N)\le C_e2^{-c(e)N}$, such that for
$G_N=\Gamma_N\setminus B_N$, all $2^{-N}\le u\le1$, and
all spherical caps $B(w,u)$,
\begin{align*}
 \mu_2\{y:(x,y)\in G_N,\ \pi_x(y)\in B(w,u)\}
     &\le2^{eN}u^S,\\
 \mu_1\{x:(x,y)\in G_N,\ \pi_y(x)\in B(w,u)\}
     &\le2^{eN}u^S.
\end{align*}
\end{enumerate}
The last bounds integrate the original opposite marginal, without
normalizing its individual fibers. The $B_N$ need not have a product
partition or be nested. All constants are independent of $N$, but
may depend on the fixed measures, $\omega$, and $e$.
\end{proposition}
\begin{proof}
Use Lemma~\ref{prep:projections} to choose an initial Frostman
probability of exponent greater than $S$. In odd dimension
$d=2k-1$, if it charges an affine $k$-plane, its normalized positive
restriction to that plane still has this exponent, and
\[
 S=k-\tfrac12\ge(k+1)/2\qquad(k\ge2).
\]
Lemma~\ref{prep:classical-plane}, applied isometrically within the
plane, proves the first alternative. These are original distances.
Otherwise project generically to $\R^{k+1}$, preserving energies
above $S$ and giving zero mass to all hyperplanes.

In both parities project further, if needed, to
$\R^{h+1}$, where $h=\lfloor S\rfloor$. All target dimensions
are greater than $S$. In even dimension this is the only auxiliary
projection. Use simultaneous potential restrictions in
Lemma~\ref{prep:projections} to obtain ball exponents above $S$
in all spaces. Choose two distinct support locations whose smallest
projected images are distinct; this is possible since that projected
law is nonatomic. Restrict to small compact positive-measure
neighborhoods of these locations. Their supports are separated in
every required space, and their original cross directions lie in a
fixed small patch. Normalize and decrease the common exponent to
$S+\chi$ if necessary.

Choose $0<\epsilon<\min(\chi,h,1)$ small enough that both laws in
$\R^{h+1}$ have finite $I_{h+\epsilon}$ energy. Apply
Lemma~\ref{prep:radial-lp} and pull its positive compact restrictions
back to the original measures. Inner regularity permits positive
compact preimages; further restriction preserves all established
bounds. These are the final $\mu_1,\mu_2$.

If $d$ is even, $h=S$. Take $\Gamma_N=E_1\times E_2$ and
apply Lemma~\ref{prep:integer-caps}, lifting to the original space.
This includes $d=2$, where the smallest projection is the identity.

If $d=2k-1$, write $\mu_i'$ for the laws in $\R^{k+1}$.
Their exponent may be decreased to an $s\in(k-\tfrac12,k)$.
They give zero mass to every affine hyperplane. Apply
Lemma~\ref{prep:coarse-filter} with the prescribed $\omega$.
The integer cap bounds obtained in $\R^k$ from
Lemma~\ref{prep:integer-caps} lift to $\R^{k+1}$, and persist
on restriction to the coarse filters. Thus $k-1$ is available there.
Proposition~\ref{prep:thin-upgrade} gives availability of
$k-\tfrac12=S$. Lift the resulting cap bounds and the exceptional
pair sets to the original space.

Pullback preserves product masses and decreasing intersections.
The auxiliary finest-depth cell partition has at most
$C2^{(k+1)\omega N}$ cells on each side. Its retained cell-pair
atoms are disjoint products. As $k+1\le d$, their number is at
most $C2^{2d\omega N}$. Pullbacks of the individual cells are
spatial Borel sets, and pullback preserves disjointness, so these
products form a partition of $\Gamma_N$. In even dimensions the
entire product is already a one-piece partition. This proves every
part of the Proposition.
\end{proof}

\begin{remark}[Order of parameters]\label{prep:parameters}
The Frostman gap $\chi$ and the measures are fixed first.
Later arguments choose $\beta,\kappa,\sigma$ using that gap, and
then choose the coarse cutoff $\omega$ sufficiently small to pay for
the product partition. Only after this choice is the plate exponent
$\zeta=\omega b$ fixed and the thin-tube upgrade invoked with its
desired loss, such as $e=\sigma$. The resulting deletion exponent
$c(e)$ may be very small. No later step requires it to dominate
$\omega$ or the Fourier saving; only its positivity is used.
\end{remark}
